% Shared exercise chunk: l2-one-bit-decisions
\begin{exercisechunk}
\item \textbf{Bayes and worst-case predictions for one bit.}
\label[exercise]{ex:one-bit-decisions}
Fix a training sample $s$ consistent with at least one hidden function, and a
test index $x$. Let $q$ be the learner's probability of predicting 1.
\begin{enumerate}[(a)]
\item If $x$ appears in $s$, show that a zero-error rule must return its
observed label with probability one.
\item Suppose $x$ is unseen and its conditional probability of label 1 is $p$
under a specified prior.
\begin{enumerate}[label=(\roman*),leftmargin=*,beginpenalty=10000]
\item Derive the conditional expected loss as a function
of $p$ and $q$.
\item Characterize its minimizing values of $q$ for every $p$.
\end{enumerate}
\item With an unknown label $b\in\{0,1\}$, write the two-by-two loss matrix
and let $r(q,b)$ be the expected loss. Solve
\[
\min_{q\in[0,1]}\max_b r(q,b),
\qquad
\max_{p\in[0,1]}\min_{q\in[0,1]}
\bigl((1-p)r(q,0)+p\,r(q,1)\bigr).
\]
Identify all optimal $q$ in the first problem and all optimal $q$ for the
least-favorable $p$ in the second. Are these sets of optimal $q$ equal?
\item Suppose an adversary chooses $b$ after seeing the learner's predicted
bit. Give an optimal strategy for this adversary, compute the resulting
minimax value, and compare it with the first value in part~(c).
\end{enumerate}
\end{exercisechunk}
